\section{Introduction}
\label{sec:intro}

Suppose you want to test a REST API that you did not write. You have its
OpenAPI specification~\cite{openapi}, which lists every operation and the
shape of its inputs, and you have the running service. Black-box testing
tools such as RESTler~\cite{restler}, RestTestGen~\cite{resttestgen} and
Schemathesis~\cite{schemathesis} take exactly these two things and send
thousands of requests, looking for server errors. REST League
2027~\cite{restleague} turns this into a competition: each tool gets one
hour per API and is scored on unique server errors (5XX responses),
operations that returned a success (2XX), and code coverage.

Most of the work in the field goes into deciding \emph{which} request to send
next: search strategies, reinforcement learning, and models of the API's
state~\cite{evomaster,arat-rl,morest,deeprest}. These methods all assume that
the tool can build a request that the API accepts. When it cannot, the
request never gets past the API's first checks, and the strategy has nothing
to work with.

We focus on that first step: turning the specification into requests that
the API accepts. It looks easy, because the specification states the type of
every input. It is not, because the API checks more than the type. A date
must parse, an id must name a resource that exists, and that resource must
have been created first. A tester that only follows the types gets a stream
of rejections, and some of them look like success: a 500 response counts as
a server error even when the only cause is a badly formed date.

Figure~\ref{fig:motivating} shows a real case from pet-clinic, one of the
REST League APIs. The operation \texttt{POST /owners/\{ownerId\}/pets} adds a
pet to an owner. A generator that follows only the schema types sent a random
string as the birth date and a random owner id. The API failed with a 500
while parsing the date, before it even looked for owner~331. This is not a
single case. Over \NReps runs, a type-only generator sent
\NMotivBaselineRequests requests to this operation; all
\NMotivBaselineServer returned 500 and none created a pet. Across the whole
API, it reached only \NPetclinicBaselineCov of the operations. In
\NSpecsWithoutExamples of the \NSpecs specifications we studied, the authors
gave no example value that a tool could copy.

\begin{figure}[t]
\begin{lstlisting}[numbers=none,escapechar=@,basicstyle=\ttfamily\scriptsize]
# Specification (resolved): body of POST /petclinic/api/owners/{ownerId}/pets
{ "name": {"type": "string"},
  "birthDate": {"type": "string", "format": "date"},  @\textbf{<- not just any string}@
  "type": {"id": {"type": "integer"}, "name": {"type": "string"}} }

# Type-only generator                     # @\tool@ (values + ids from earlier responses)
POST .../owners/331/pets                  POST .../owners/11/pets
{"name": "jzPdeIgx",                      {"name": "Leo",
 "birthDate": "LdGncfBA",                  "birthDate": "2010-09-07",
 "type": {"id": 970, "name": "epfJBdKh"}}  "type": {"id": 1, "name": "cat"}}
-> 500 "JSON parse error: Cannot          -> 201 Created
   deserialize value of type                 {"id": 14, "ownerId": 11, ...}
   java.time.LocalDate from String
   'LdGncfBA' ..."
\end{lstlisting}
\caption{\textcolor{red}{[Motivating example --- graded]} A real exchange with
  pet-clinic. Following only the schema types, a generator sends a random
  birth date, and the API fails while parsing it; over \NReps runs, all
  \NMotivBaselineRequests such requests returned 500 and none created a pet.
  \tool sends a real date, an existing owner and an existing pet type, and
  the pet is created. Requests and responses are copied from
  \texttt{data/motivating.json}, trimmed with \texttt{...}.}
\label{fig:motivating}
\end{figure}

Earlier work has attacked parts of this problem. ARTE~\cite{arte} and
NLP2REST~\cite{nlp2rest} look outside the schema, in knowledge bases or in
the text of the documentation, to find realistic values. RESTler and
RestTestGen infer which operation creates the resource another needs.
However, these parts are built into complete tools, so it is hard to tell
how much each one adds. A competitor who has a few weeks must decide what to
build first.

We present \tool, the part of our REST League tool that turns a
specification into requests. It works in four steps:
\begin{enumerate}
  \item it parses the specification and resolves every reference, merged
    schema and recursive schema (Section~\ref{sec:step1});
  \item it orders the operations so that those that create a resource come
    before those that use it (Section~\ref{sec:step2});
  \item it fills each request with values taken from examples, formats,
    regular-expression patterns and parameter names, in a valid mode for
    reaching operations and two error modes for finding faults
    (Section~\ref{sec:step3});
  \item it reuses real ids from earlier responses through the engine's
    resource pool (Section~\ref{sec:step4}).
\end{enumerate}
\tool uses only the Python standard library, no network and no language
model, so it fits the competition's sandbox.

We evaluated \tool on all \NSpecs specifications of the RESTgym
APIs~\cite{restgym} and on two REST League APIs running live. \tool parsed
every specification with no reference left unresolved, and \NConfTool of its
generated values followed their schema, against \NConfBaseline for a
type-only generator. With the same budget, the operations reached rose from
\NPetclinicBaselineCov to \NPetclinicFullCov on pet-clinic and from
\NGestaohospitalBaselineCov to \NGestaohospitalFullCov on gestao-hospital.
Taking each part away in turn showed where the gain comes from: values and
id reuse carried it, while operation order added no reached operation on
either API. The lesson for a team with little time is to get the values and
the ids right first.

\paragraph{Contributions.}
\begin{itemize}
  \item We state the input side of black-box REST testing as reaching
    operations under a request budget (Section~\ref{sec:problem}).
  \item We present \tool, a small, dependency-free component that parses
    OpenAPI specifications, generates values in three modes, reuses ids and
    orders operations (Section~\ref{sec:approach}).
  \item We measure each part on its own, on \NSpecs specifications and two
    live APIs, and find that values and id reuse matter most
    (Section~\ref{sec:eval}).
  \item We release the code, the scripts and every raw result
    (Section~\ref{sec:data}).
\end{itemize}
